\subsubsection*{1.6.}
沿用 1.1 的假设. 设 $\bar s$ 为 $S$ 的几何点, 是 $\bar x_0$ 的像; 设 $S_{(\bar s)}$ 为 $\mathcal O_{S,s}$ 在 $\bar s$ 处的严格 Hensel 化的谱, $X_{(\bar s)}$ 为 $X$ 在 $S_{(\bar s)}$ 上的逆像. 设 $\bar\eta$ 为 $S_{(\bar s)}$ 的分式域 $k(\eta)$ 的一个代数闭包的谱, 置
\[
I=\operatorname{Gal}(k(\bar\eta)/k(\eta)).
\]

设 $\ell$ 为不同于 $S$ 的剩余特征的素数. 在 $X_{\bar s}$ 中 $\bar x_0$ 的一个邻域内, 关于 $X_{(\bar s)}$ 上常值层 $\mathbb Z/\ell$ 的消灭圈层
\[
\Phi^i_{\bar\eta}(\mathbb Z/\ell)
=R^i\Phi_{\bar\eta}(\mathbb Z/\ell)
\]
集中在 $\bar x_0$ 处. 此外, 惯性群 $I$ 作用于 $\mathbb Z/\ell$-向量空间
\[
\Phi^i_{\bar\eta}(\mathbb Z/\ell)_{\bar x_0}.
\]
这些空间是有限维的（XIII, 2.4.2）.

\subsubsection*{1.7.}
设 $I_1$ 为惯性群的一个有限商群. 回顾 $I_1$ 的 \underline{Swan 特征标} $\operatorname{Sw}$ 是 Artin 特征标 $a$ 与增广表示的特征标 $u$ 之差 $a-u$. 还要回顾, 在同构意义下存在唯一的射影 $\mathbb Z_\ell[I_1]$-模 $S(I_1)$, 其特征标为 $\operatorname{Sw}$（见 [6]）.

设 $V$ 为有限维 $\mathbb Z/\ell$-向量空间, 带有作用 $\rho:I\to\operatorname{GL}(V)$, 并假设 $\rho$ 通过 $I_1$ 分解. 则 $V$ 的\underline{野维数}（wild dimension）为
\begin{equation}\tag{1.7.1}
\dim\operatorname{Sw}(V)
=\dim_{\mathbb Z/\ell}
\operatorname{Hom}_{I_1}\!\left(
S(I_1)\otimes_{\mathbb Z_\ell}\mathbb Z/\ell,,V\right).
\end{equation}
此整数不依赖于 $I_1$ 的选择, 因为对 $I_1$ 的正规子群 $H$, 有 $S(I_1/H)\sim S(I_1)_H$.

再置
\[
\dim\operatorname{tot}(V)
=\dim\operatorname{Sw}(V)+\dim_{\mathbb Z/\ell}(V)
\]
称此数为 $V$ 的\underline{总维数}（total dimension）.

\subsubsection*{定义 1.8.}
\underline{采用 1.6 与 1.7 的术语, $X/S$ 在 $\bar x_0$ 处的消灭圈数（分别为消灭圈的野数、消灭圈的总数）为下列整数}
\[
(-1)^n\dim\Phi_{\bar x_0}
\underset{\mathrm{def}}{=}
(-1)^n\sum_i(-1)^i
\dim\Phi^i_{\bar\eta}(\mathbb Z/\ell)_{\bar x_0},
\]
\[
(-1)^n\dim\operatorname{Sw}\Phi_{\bar x_0}
\underset{\mathrm{def}}{=}
(-1)^n\sum_i(-1)^i
\dim\operatorname{Sw}
\Phi^i_{\bar\eta}(\mathbb Z/\ell)_{\bar x_0},
\]
\[
\begin{aligned}
(-1)^n\dim\operatorname{tot}\Phi_{\bar x_0}
&\underset{\mathrm{def}}{=}
(-1)^n\sum_i(-1)^i\\[-2pt]
&\qquad{}\cdot\dim\operatorname{tot}
\Phi^i_{\bar\eta}(\mathbb Z/\ell)_{\bar x_0}\\
&=(-1)^n\bigl(
\dim\Phi_{\bar x_0}+\dim\operatorname{Sw}\Phi_{\bar x_0}
\bigr).
\end{aligned}
\]

\subsubsection*{猜想 1.9.}
\underline{在 1.1 的假设下, $X/S$ 在 $\bar x_0$ 处的 Milnor 数 $\mu$ 等于 $X/S$ 在 $\bar x_0$ 处的消灭圈总数.}

\subsubsection*{注 1.10.}
设 $S$ 为剩余域不完美的 Hensel 离散赋值环谱, $f:X\to S$ 为除 $x_0\in X_s$ 外处处光滑的态射. 假设 $X$ 在 $x_0$ 处正则. 在这些假设下, Milnor 数 $\mu$ 的定义 1.2 仍有意义, 但我已无法为 $\mu$ 提出几何解释.
